\documentclass{article}
\usepackage{geometry}
\usepackage{concmath}
\usepackage[OT1]{fontenc}
\usepackage{graphicx}
\usepackage{amsfonts}
\usepackage{amsmath,mathtools}
\usepackage{calrsfs}
\usepackage{color}
\usepackage[dvipsnames]{xcolor}
\usepackage[many]{tcolorbox}
\usepackage{setspace}
\usepackage{tkz-tab}    
\usepackage{amssymb}          
\usepackage{multicol}  
\geometry{margin=0.5in}
\definecolor{green}{RGB}{52, 235, 168}
\definecolor{violet}{RGB}{165, 83, 237}
\definecolor{lemon}{RGB}{170, 237, 83}
\definecolor{red}{RGB}{240, 125, 84}
\definecolor{blue}{RGB}{105, 255, 242}

\begin{document}
\noindent \colorbox{green}{640.2}
\begin{align*}
    I_{n+2} &= \int_{0}^{\frac{\pi}{2}} sint sin^{n+1}t dt \\
    &= [-cost sin^{n+1}t]^{\frac{\pi}{2}}_0 - \int_{0}^{\frac{\pi}{2}}(-cost)(n+1)cost sin^nt dt \\
    &= 0 + (n+1)\int_{0}^{\frac{\pi}{2}} (1-sin^2t)sin^nt dt \\
    &= (n+1)(I_n - I_{n+2})
\end{align*}
donc \(I_{n+2} = \frac{n+1}{n+2}I_n\)
\begin{align*}
    I_{2p} = \frac{(2p)!}{(2^p p!)^2}\frac{\pi}{2} \indent \text{ hérédité}
\end{align*}
Au voisinage de \(p \to +\infty\)\\
\begin{align*}
    I_{2p} &= \frac{(2p)!}{2^{2p}(p!)^2}\frac{\pi}{2}\\
    &\sim \frac{\sqrt{2\pi}\sqrt{2p}(2p)^{2p}e^{-2p}}{2^{2p}2\pi p p^{2p}e^{-2p}}\frac{\pi}{2}\\
    &= \sqrt{\pi}\frac{1}{2}\frac{1}{\sqrt{p}} 1\\
    &= \sqrt{\frac{\pi}{4p}}
    \rightarrow_{p \to +\infty}0
\end{align*}

\noindent \colorbox{green}{640.3} \begin{align*}
    u_n &= \sum_{k=1}^{n}\frac{1}{n^2+k^2}\\
    &= \frac{1}{n}(\frac{1}{n}\sum_{k=1}^{n}\frac{1}{(\frac{k}{n})^2 + 1})
\end{align*}
On reconnaît une somme de Riemann avec la fonction \(f : x \mapsto \frac{1}{1+x^2}\)\\
Sur le segment [0,1]:\\
* f est continue sur [0,1]\\
donc \begin{align*}
    \frac{1}{n}\sum_{k=1}^{n}f(\frac{k}{n}) &\rightarrow_{n\to +\infty} \int_{0}^{1}f(t)dt\\
    &= Arctan 1\\
    &= \frac{\pi}{4}
\end{align*}
Ainsi \begin{equation*}
    u_n \sim_{n\to +\infty}\frac{1}{n}\frac{\pi}{4}
\end{equation*}
\colorbox{green}{640.5}\\
a,b \(\in \mathbb{R}\) tq \(a < b\)\\
h continue positive\\
On veut montrer \begin{equation*}
    \int_{0}^{1}h(x)dx = 0 \Rightarrow h = 0
\end{equation*}
Par contraposée on suppose \(h \neq 0\)\\
\begin{equation*}
    \exists x_0 \in [a,b], h(x_0) > 0
\end{equation*}
Par définition \(\varepsilon = \frac{h(x_0)}{2}\)\begin{align*}
    &\exists \alpha > 0, x \in [x_0-\alpha, x_0 + \alpha] \Rightarrow |h(x)-h(x_0)| \leq \varepsilon\\
    & \forall x \in [x_0-\alpha, x_0 + \alpha], h(x)>\varepsilon > 0\\
    \int_{0}^{1} h(x)dx &= \int_{0}^{x_0-\alpha} h(x) dx + \int_{x_0-\alpha}^{x_0+\alpha} h(x) dx + \int_{x_0+\alpha}^{1} h(x)\\
    &\geq 0 + 2\alpha \frac{h(x_0)}{2} + 0\\
    &\geq 0
\end{align*}
\colorbox{green}{640.17}\\
\begin{equation*}
    \frac{1}{2}ln((t-a)^2+b^2) + iArctan(\frac{t-a}{b})
\end{equation*}

\noindent \colorbox{green}{640.4}\\
\begin{equation*}
    sin^{(k)}(x) = sin(x + k\frac{\pi}{2})
\end{equation*}
Soit \(x \in \mathbb{R}\) fixé : \begin{align*}
    sinx &= \sum_{k=0}^{n}sin^{(k)}(0)\frac{x^k}{k!} + \int_{0}^{x}sin^{(n+1)}(t) \frac{(x-t)^n}{n!}dt\\
    &\int_{0}^{x}sin^{(n+1)}(t) \frac{(x-t)^n}{n!}dt = R_n(x)
\end{align*}
Mq \(R_n(x)\rightarrow_{n \to +\infty} 0, x>0\)\begin{align*}
    |R_n(x)| &\leq \int_{0}^{x}|sin^{(n+1)}(t)\frac{(x-t)^n}{n!}|dt\\
    &\leq \int_{0}^{x}1 \frac{(x-t)^n}{n!} dt\\
    &= \frac{1}{n!}[\frac{(x-t)^{n+1}}{(n+1)!}]^x_0\\
    &= \frac{x^n+1}{(n+1)!}\\
    &\rightarrow_{x \to +\infty} 0
\end{align*}
Donc \begin{equation*}
    sinx = \sum_{k=0}^{+\infty} sin^{(k)}(0)\frac{x^k}{k!}
\end{equation*}




\end{document}